\section{Conclusion}\label{sec:conclusion}

GhostVault keeps files in private GitHub repositories without trusting GitHub with them. The
cryptography is conventional: a random key per file, AES-256-GCM per chunk, associated data that binds
each chunk to its file, position, count, size and padding, and keys for recovery bundles derived with
Argon2id. Most of the system is about something else: deleting the user's data only when another copy
has been shown to exist. The original goes only after the branch has been read back and a sample
downloaded, and after the file has been moved aside and checked again; the remote copy goes only after
the restored file is on disk under its final name. Every step that changes something outside the
computer is journaled first, every reply that matters is checked against what the repository shows,
and every record is validated before anything acts on it.

The protection has clear edges. Cryptography protects the content of files from everyone who can read
the repositories; GitHub's access control decides who those readers are; and only GitHub keeps the data
available, which is why a ghost should never be the only copy of a file that cannot be lost. Within
those edges, the system was checked by an automated suite covering every journal checkpoint, by
randomised runs in which 1\,273 simulated crashes and 7\,736 network faults lost no file, and by
measurements showing that the limit on its speed is GitHub's pacing, not its own work.

\section*{Availability}
GhostVault is open source under the MIT license. The source, its documentation, the test suite, the
measurement and stress tools, and this report's \LaTeX{} source and data are in the project's
repository.
